\subsection*{4.2.1\quad 计算芯片：为什么是 GPU？}

先看大模型到底在算什么。无论是注意力中的投影和加权求和，还是前馈网络中的两次线性变换，Transformer 的绝大部分计算都可以归结为矩阵乘法。矩阵乘法有一个很好的性质：它由海量彼此独立的乘加运算组成，谁先算、谁后算都不影响结果。这意味着，决定速度的不是单个运算有多快，而是\emph{同一时刻能有多少运算并行进行}。

CPU 恰好是按相反的思路设计的。它只有几十个核心，但每个核心都配有复杂的分支预测和多级缓存，目标是让一条指令流尽快执行完毕——这对操作系统、网络服务这类控制流密集的任务非常合适，对矩阵乘法却是“大材小用”。GPU 则把芯片面积花在成千上万个简单的执行单元上，并配备专门的矩阵乘加单元（如 NVIDIA 的 Tensor Core）和高带宽显存，单个线程很慢，但整体吞吐极高。

沿着“为特定计算专门设计”的方向再往前走，就得到 NPU 与 ASIC。这两个名称经常并列出现，但含义并不相同：NPU（神经网络处理器）是按\emph{用途}命名的，ASIC（专用集成电路）是按\emph{实现方式}命名的，一款 NPU 完全可以用 ASIC 的方式实现。Google TPU、华为昇腾和寒武纪 MLU 都属于这一类（图 4-3）。

这样看来，三类芯片并不是简单的“谁更好”，而是在通用性与专用效率之间各取一点。实际系统中它们也是协作关系：CPU 负责调度、预处理和下发任务，GPU 或 NPU 负责吞吐密集的张量运算。

专用化带来了更高的吞吐和能效，代价是什么？一是通用性下降，二是软件生态的绑定更深。这也提醒我们，比较芯片时不能只看厂商给出的“TFLOPS”。同一块芯片在 FP16、FP8 或 INT8 下的峰值可以相差数倍，是否计入稀疏加速也会让数字翻倍。真正有意义的比较，要同时看目标模型能否运行、支持哪些低精度格式、显存有多大、编译器和算子库是否成熟，以及集群通信能力和总拥有成本。

\subsection*{4.2.2\quad 存储层次：为什么 HBM 如此重要}

第 4.1.2 小节留下了一个结论：解码时，每生成一个 token 都要把全部权重从显存读一遍，速度上限由显存带宽决定。训练时情况更紧张——权重、梯度、优化器状态和激活值都要放在设备上。于是存储系统要同时回答两个问题：\emph{装不装得下}，以及\emph{送不送得及}。

单一一种存储无法同时做到又大、又快、又便宜，计算机的办法是分层（图 4-4）。离计算单元越近，存储越快、越小、越贵：片上 SRAM 带宽最高，但只有几十 MB；板载显存有几十 GB；主机内存可达 TB 级，可 GPU 要经过 PCIe 才能访问它，带宽低了一个数量级以上；SSD 则用来存放模型文件和冷数据。

在这个层次中，显存这一级对大模型最为关键，而当前数据中心加速器的主流显存方案是\term{高带宽内存}（HBM）。它的思路很直接：既然带宽受限于接口宽度和传输距离，那就把多层 DRAM 芯片垂直堆叠起来，用硅通孔（TSV）连接各层，再以每个堆栈 1024 位的超宽接口与计算芯片封装在一起。接口变宽、距离变短，带宽自然大幅提高。以 H100 SXM 为例，它配备 \SI{80}{GB} HBM3，带宽约 \SI{3.35}{TB/s}。

不过，HBM 解决的主要是“送不送得及”，而不是“装不装得下”。它价格昂贵，容量依然有限；运行大模型也并非必须使用 HBM——消费级显卡的 GDDR 显存同样可以运行规模适中或经过量化的模型。不妨把三个指标的分工记在一起：\emph{容量决定能否装下，带宽决定能否及时送达，延迟影响细粒度访问的效率}。

\subsection*{4.2.3\quad 互联：数据要从哪里传到哪里？}

一张卡总有装不下、算不完的时候。例 4-1 估算过，用一张卡训练一个 70B 模型要花约 6.7 年，想在几天内完成就必须让上千张卡并行；即使只做推理，较大的模型也常常要切分到多张卡上。一旦切分，卡与卡之间就要不断交换梯度、激活值和中间结果。此时再快的计算单元，如果数据迟迟不到，也只能空等。

讨论互联时，最有用的问题不是“哪种技术更快”，而是“\emph{数据要从哪里传到哪里}”。不同层级的链路各有分工（图 4-5）：
\begin{itemize}
  \item 在一台服务器内部，CPU 与 GPU、网卡、硬盘之间通过通用总线 \term{PCIe} 相连，PCIe 5.0 x16 的单向带宽约 \SI{64}{GB/s}；
  \item 同一台服务器内的 GPU 之间，NVIDIA 用专门的 \term{NVLink} 直连，H100 每卡总带宽约 \SI{900}{GB/s}，再借助 NVSwitch 让 8 张卡两两全速互通；
  \item 跨越服务器时，就要靠 \term{InfiniBand} 这样的集群网络，单端口可达 \SI{400}{Gb/s}，并原生支持 RDMA——数据可以直接写入远端机器的内存，省去双方 CPU 的多次拷贝。
\end{itemize}

把这三个数字放在一起，就能看出一条贯穿全章的规律：链路带宽逐级相差一个数量级，因此\emph{通信越频繁的工作，越应该放在越近的地方}。第 4.3.4 小节会看到，张量并行每一层都要通信，应当留在 NVLink 相连的同一台机器内；数据并行和流水线并行通信较少，才适合跨越 InfiniBand。
